% Circuito RC em série
% Feito para o ÉguaTeX (eguatex.com.br). Uso livre (CC0 1.0).
%
% Compile e pronto: sai só a figura, em PDF.
% Para pôr no seu trabalho, use o PDF com \includegraphics ou copie
% o bloco circuitikz para dentro de uma figure e ponha no preâmbulo
% \usepackage[american]{circuitikz}.
% Para os símbolos europeus (resistor em retângulo), troque american por european.
\documentclass[border=4pt]{standalone}
\usepackage[T1]{fontenc}
\usepackage{lmodern}
\usepackage[american]{circuitikz}

\begin{document}
\begin{circuitikz}[line cap=round, line join=round]

% Mude estes textos e tamanhos.
\def\fonte{$\varepsilon$}   % fonte, por exemplo $12\,\mathrm{V}$
\def\resistor{$R$}         % resistor, por exemplo $1\,\mathrm{k\Omega}$
\def\capacitor{$C$}        % capacitor, por exemplo $100\,\mu\mathrm{F}$
\def\largura{5}            % largura do circuito, em cm
\def\altura{3.5}           % altura do circuito, em cm

\draw (0,0) to[battery1, l=\fonte, invert] (0,\altura)
            to[spst, l=$S$] ({\largura/2},\altura)
            to[R, l=\resistor, i>^=$i$] (\largura,\altura)
            to[C, l_=\capacitor, v^>=$v_C$] (\largura,0)
            -- (0,0);
\end{circuitikz}
\end{document}
